\centerline{\Huge \bf  Training Olympiad 6}
\centerline{\Huge \bf Level: IMO (IMC)}

\begin{flushright}
    \scriptsize{Chapter 6. Knockin' on Heaven's Door}
\end{flushright}

\noindent\textbf{Problem 1.} We colour all the sides and diagonals of a regular polygon P with 43 vertices either red or blue in such a way that every vertex is an endpoint of 20 red segments and 22 blue segments. A triangle formed by vertices of P is called monochromatic if all of its sides have the same colour. Suppose that there are 2022 blue monochromatic triangles. How many red monochromatic triangles are there?

\noindent\textbf{Problem 2.} Let $p > 2$ be a prime number. Prove that there is a permutation $(x_1, x_2, \ldots, x_{p-1})$ of the numbers ($1, 2, ..., p - 1)$ such that:
\[x_1x_2 + x_2x_3 + \ldots + x_{p-2}x_{p-1} \equiv 2 \ (mod \ p).\]

\noindent\textbf{Problem 3.} Let $ABC$ be an acute-angled triangle and let $D, E,$ and $F$ be the feet of altitudes from $A, B,$ and $C$ to sides $BC, CA,$ and AB, respectively. Denote by $\omega_{B}$ and $\omega_{C}$ the incircles of triangles $BDF$ and $CDE$, and let these circles be tangent to segments DF and DE at M and N, respectively. Let line MN meet circles $\omega_{B}$ and $\omega_{C}$ again at $P \neq M$ and $Q \neq N$, respectively. Prove that $MP$ = $NQ$.